% !TeX program = xelatex
\documentclass[12pt,a4paper]{ctexart}

\usepackage{amsmath,amssymb}
\usepackage{booktabs}
\usepackage{geometry}
\usepackage{hyperref}
\usepackage{enumitem}
\usepackage{array}

\geometry{left=2.5cm,right=2.5cm,top=2.5cm,bottom=2.5cm}
\setlength{\parskip}{0.4em}
\hypersetup{colorlinks=true,linkcolor=blue}

\title{\textbf{数学建模项目开题报告}\\[0.4em]
\large 恐惧效应下捕食者—猎物动力学建模与数值分析}
\author{%
  李松茂（PB23151824）\quad
  贺小轩（PB23151820）\quad
  王玉麟（PB23151808）%
}
\date{2026 年 5 月 28 日}

\begin{document}
\maketitle

\begin{center}
\small
\textbf{所属课程：}数学建模 \quad
\textbf{项目类型：} predator--prey dynamics with fear \& memory effects
\end{center}

\vspace{0.5em}
\hrule
\vspace{1em}

\section{选题背景与意义}

\subsection{选题来源}

本项目选题来自课程大作业里的\textbf{「恐惧效应下的捕食者—猎物动力学」}：在经典种群模型中引入\textbf{恐惧因子}或\textbf{记忆效应}，研究非消耗性效应（NCE）对共存、振荡与崩溃的影响，并与无恐惧基线对比，通过参数实验与敏感性分析说明定性机制。

\subsection{选题依据}

经典 Lotka--Volterra 与 Rosenzweig--MacArthur 模型主要刻画捕食者\textbf{直接消耗}猎物（CE）。近二十年来大量实验与综述表明：猎物在\textbf{未被捕食}时也会因感知风险而改变行为与繁殖，其效应强度常与 CE 同量级（Preisser et al., 2007; Clinchy \& Sheriff, 2011）。Zanette et al.\ (2003) 在 \emph{PNAS} 报道捕食者叫声可使繁殖成功率下降约 53\%；Peacor \& Werner (2001) 证明仅视觉线索即可改变水蚤种群动态。因此，在数学建模竞赛框架下，将\textbf{恐惧 + 记忆}写入 ODE，并结合\textbf{公开生态数据}与\textbf{文献机制}交叉验证，具有明确的科学问题与可操作的数值路线。

\subsection{研究意义}

\begin{itemize}[leftmargin=2em]
  \item \textbf{理论意义：}连接 Wang et al.\ (2019)「恐惧与稳定性」与 Myint et al.\ (2025) B--D+恐惧拓展，回应「恐惧能否稳定化周期系统」；
  \item \textbf{方法意义：}建立「机制建模—参数扫描—数据标定—文献对照」的分析框架；
  \item \textbf{应用意义：}为陆生哺乳类、淡水鱼类、浮游动物等多类捕食—猎物系统提供统一的恐惧效应比较视角。
\end{itemize}

\section{研究思路与技术路线}

\subsection{总体思路}

采用\textbf{双主线 + 多层验证}：

\begin{enumerate}[label=\textbf{(\arabic*)}]
  \item \textbf{作业主线（恐惧+记忆 ODE）：}在逻辑斯蒂--Holling\,II 基线上建立繁殖抑制型恐惧项与指数核记忆变量 $M(t)$，满足题目对「记忆效应」的明确要求；
  \item \textbf{拓展主线（B--D+恐惧）：}参照 Myint (2025)，用 Beddington--DeAngelis 功能反应与饱和因子 $f(k,v)=1/(1+kv)$ 对接真实时间序列；
  \item \textbf{机制对照：}按文献实现繁殖抑制、觅食抑制、处理时间延长等多种 NCE 刻画，说明「恐惧进入方程的路径」决定定性结论（Abrams, 2000）；
  \item \textbf{恐惧稳定化检验：}以平衡点 Jacobian 特征值、长期振幅与峰值衰减比为框架，检验恐惧参数对共存与振荡的调节；
  \item \textbf{数据与文献验证（A/B 双轨）：}\textbf{A 轨}——12 条序列标定后跨类群比较等效抑制 $\eta(k,v)$ 及组内异质（Andrén/Killifish）；\textbf{B 轨}——Peacor 元分析、珊瑚/豆娘实验、LTER 追加拟合作为独立于标定的机制先验；两轨互补、不可混读。
\end{enumerate}

\subsection{技术路线（流程）}

\begin{center}
\fbox{\parbox{0.92\textwidth}{\small
\textbf{文献与数据调研} $\rightarrow$ \textbf{主模型 + B--D 拓展建立} $\rightarrow$
\textbf{数值实验}（$\phi,\delta,k$ 扫描、机制对照、稳定化分析）\\
$\rightarrow$ \textbf{A 轨 12 序列标定 + B 轨实验/元分析先验（双轨验证）}\\
$\rightarrow$ \textbf{可选 2D PDE 拓展 + Turing 线性诊断（附录，不作主结论）}\\
$\rightarrow$ \textbf{论文撰写与答辩}
}}
\end{center}

\subsection{核心模型（摘要级表述）}

\textbf{主模型（恐惧+记忆）：}
\begin{equation*}
\frac{dx}{dt} = rx\left(1-\frac{x}{K}\right) - \frac{axy}{1+\theta x} - r\phi M x,\quad
\frac{dy}{dt} = \frac{eaxy}{1+\theta x} - \mu y,\quad
\frac{dM}{dt} = y - \delta M.
\end{equation*}

\textbf{拓展模型（B--D+恐惧，用于数据标定）：}
\begin{equation*}
\frac{du}{dt} = \frac{ru}{1+kv} - du - au^2 - \frac{puv}{1+qu+v},\quad
\frac{dv}{dt} = v\left(-m + \frac{cpu}{1+qu+v}\right).
\end{equation*}

\section{总体规划与进度安排}

\subsection{阶段划分}

\begin{table}[h]
  \centering
  \caption{项目四阶段规划}
  \small
  \begin{tabular}{p{2.2cm}p{5.2cm}p{6.2cm}}
    \toprule
    阶段 & 主要任务 & 预期产出 \\
    \midrule
    第一阶段 & 文献调研、问题重述、模型假设与方程推导 & 开题报告、符号表、理论框架 \\
    第二阶段 & 数值程序实现、$\phi/\delta/k$ 扫描、机制对照与稳定化分析 & 主实验数值结果、参数扫描图 \\
    第三阶段 & 公开数据整理、多模型标定、机制先验对照 & 拟合曲线、参数汇总表、对照分析 \\
    第四阶段 & 论文统稿、图表整理、答辩 PPT & 竞赛论文、答辩材料 \\
    \bottomrule
  \end{tabular}
\end{table}

\subsection{工作与成果结构（规划）}

开题阶段数值实验与标定尚未开展，本节仅规划后续工作的大致模块划分：

\begin{itemize}[leftmargin=2em]
  \item \textbf{建模与数值模块：}ODE 右端项（基线、恐惧+记忆、B--D+恐惧及机制对照）、数值积分、参数扫描与敏感性分析；可选 2D 反应—扩散拓展；
  \item \textbf{数据与标定模块：}公开数据下载与元数据整理、序列识别、多模型拟合与参数估计；
  \item \textbf{验证与对照模块：}将模型参数与文献元分析、行为实验及野外监测结论进行机制层面对照；
  \item \textbf{文档与答辩模块：}竞赛论文、进度备忘、复现说明及答辩演示材料。
\end{itemize}

\section{调研成果}

\subsection{文献调研结论}

\begin{enumerate}[label=\arabic*.]
  \item \textbf{NCE 与 CE 同等重要：}Preisser (2007) 元分析表明性状介导效应（TMI）与密度介导效应（DMI）强度相当，模型须显式区分「仅捕食」与「捕食+恐惧」；
  \item \textbf{恐惧可大幅压低繁殖：}Zanette (2003) 实验支持将恐惧写入有效繁殖率 $r_{\text{eff}}$ 或等价抑制因子；
  \item \textbf{记忆可用 ODE 化：}MacDonald (1989) 指数核 $K(\tau)=\delta e^{-\delta\tau}$ 可化为 $dM/dt=y-\delta M$，便于与经典 ODE 框架衔接；
  \item \textbf{机制路径敏感：}Abrams (2000) 强调不同 NCE 通道（繁殖抑制、觅食抑制、干扰竞争等）可给出相反符号结论，须做多机制对照而非单一恐惧项；
  \item \textbf{恐惧与稳定性：}Wang (2019)、Myint (2025) 指出恐惧参数可能改变平衡点位置、Jacobian 谱及极限环形态——后续数值实验将围绕该命题展开。
\end{enumerate}

\subsection{数据调研与拟用来源}

开题阶段已完成公开数据源的初步检索，拟从以下来源中选取具有明确捕食—猎物结构的长时间序列用于标定与对照（具体条数与可用性将在第二阶段核实）：

\begin{table}[h]
  \centering
  \caption{拟调研/选用的主要数据来源}
  \small
  \begin{tabular}{llp{5.5cm}}
    \toprule
    平台/文献 & 典型系统 & 调研关注点 \\
    \midrule
    Dryad（Andrén 等） & 欧亚猞猁—狍 & 多区域、陆生哺乳类 \\
    Dryad（Marino 等） & 雪兔—猞猁 & 经典周期种群 \\
    GLERL & 密歇根湖浮游 & 入侵—浮游相互作用 \\
    Killifish 数据集 & 鱼类—食蚊鱼 & 多站点、淡水系统 \\
    Peacor (2022) & 捕食风险效应文献清单 & TMIE/NCE 机制先验 \\
    LTER / 行为实验数据 & 鱼类丰度、活动度等 & 独立机制验证 \\
    \bottomrule
  \end{tabular}
\end{table}

\subsection{调研阶段小结}

基于上述文献与数据检索，开题阶段形成以下\textbf{待验证}认识（尚未开展数值实验，不作为最终结论）：

\begin{itemize}[leftmargin=2em]
  \item 恐惧效应不宜仅用单一参数 $\phi$ 概括，须区分进入繁殖项、功能反应或干扰项的不同通道（Abrams, 2000）；
  \item 记忆效应可通过指数核引入，与题目「记忆」要求一致，且 $\delta$ 可解释为风险信息遗忘速率；
  \item B--D+恐惧模型（Myint, 2025）在结构上兼顾捕食干扰与饱和恐惧，适合作为对接真实序列的拓展主线；
  \item Peacor 等元分析提示 TMIE 在文献中占相当比例，支持在本项目中保留非消耗性效应通道；
  \item 跨系统标定时，恐惧参数的量纲与可识别性可能因数据质量而异，拟采用等效抑制量 $\eta(k,v)$ 辅助解读，并采用\textbf{A/B 双轨}分别回答「模型可迁移性」与「NCE 机制是否有独立证据」。
\end{itemize}

\subsection{阶段性成果（与竞赛论文同步，2026 年 5 月）}

\begin{itemize}[leftmargin=2em]
  \item \textbf{12/12 序列} B--D+恐惧 RMSE 最优（0.054--0.209），较 baseline 改进 $10^2$--$10^6$ 倍；
  \item \textbf{A 轨：}哺乳/鱼类 $\eta(v=5)$ 中位数 $\sim$0.07--0.08；Andrén/Killifish 组内异质大于类群间差异；
  \item \textbf{B 轨：}Peacor 64 篇 TMIE 58\%；无脊椎猎物 TMIE 11/12（92\%）；豆娘活动度 $\sim$6\% 与 A 轨 $\eta$ 同阶；
  \item \textbf{2D 拓展：}已实现有限差分 + $d_2$--$\max_k\mathrm{Re}\,\lambda(k)$ 线性诊断；PDE 未见显著 Turing 斑图，\textbf{正文不作主结论}（拓展性尝试，非实验失败）。
\end{itemize}

\section{组员分工}

\begin{table}[h]
  \centering
  \caption{三人分工表}
  \begin{tabular}{p{2.8cm}p{2.5cm}p{7.5cm}}
    \toprule
    姓名（学号） & 角色定位 & 主要职责与交付物 \\
    \midrule
    贺小轩（PB23151820） &
    模型与理论 &
    问题重述与假设；从经典 LV 到恐惧—记忆 ODE 及 B--D 拓展的公式推导；平衡点与稳定性分析；恐惧稳定化检验框架设计；参数与机制的理论解读；论文模型建立与理论分析章节。 \\
    \midrule
    李松茂（PB23151824） &
    编程与数值实验 &
    数值程序设计与实现；参数扫描、机制对照与稳定化分析实验；公开数据标定流程搭建；实验结果整理与可复现性保障；配合模型调试与答辩数值演示。 \\
    \midrule
    王玉麟（PB23151808） &
    文献、数据与论文 &
    文献检索与综述（Preisser、Wang、Myint、Peacor 等）；公开数据调研、下载与引用规范；论文摘要、引言、结论、参考文献及全文统稿；图表排版与 PDF 编译；开题/答辩 PPT 与汇报材料。 \\
    \bottomrule
  \end{tabular}
\end{table}

\subsection{协作机制}

协作安排与上述四阶段规划对应，各阶段内按「模型—编程—论文」衔接：

\begin{itemize}[leftmargin=2em]
  \item \textbf{第一阶段（文献与建模）：}王玉麟整理文献与数据清单 $\rightarrow$ 贺小轩完成问题重述、假设与方程推导、符号表 $\rightarrow$ 三人确认理论框架与开题报告；
  \item \textbf{第二阶段（数值实验）：}李松茂实现程序并完成 $\phi/\delta/k$ 扫描、机制对照与稳定化分析 $\rightarrow$ 贺小轩解读数值结果与稳定性 $\rightarrow$ 王玉麟记录实验要点、预备论文数值节素材；
  \item \textbf{第三阶段（数据标定与对照）：}王玉麟协助公开数据整理与引用规范 $\rightarrow$ 李松茂开展多模型标定与机制先验对照 $\rightarrow$ 贺小轩从理论角度解读参数与跨系统差异 $\rightarrow$ 王玉麟写入标定与对照分析；
  \item \textbf{第四阶段（论文与答辩）：}三人交叉审阅竞赛论文；王玉麟统稿与图表排版，贺小轩把关理论表述，李松茂负责数值部分演示与答辩答疑准备。
\end{itemize}

\section{预期成果与创新点}

\subsection{预期成果}

\begin{enumerate}[label=\arabic*.]
  \item 完整的捕食—猎物恐惧效应数值分析程序及实验记录；
  \item 数学建模竞赛论文（含机制建立、数值分析、数据标定与模型评价）；
  \item 开题/答辩材料（本报告、PPT、复现说明文档）。
\end{enumerate}

\subsection{拟创新点}

\begin{itemize}[leftmargin=2em]
  \item \textbf{A/B 双轨恐惧验证：}A 轨跨类等效 $\eta$ 比较 + B 轨 Peacor/实验/元分析独立先验，相互印证但不可混读；
  \item \textbf{机制—数据—稳定化}三位一体：在参数扫描之外，对接 12 序列标定与 Peacor/LTER 等独立机制先验；
  \item \textbf{双主线模型：}记忆 ODE 满足题目要求，B--D+恐惧 服务真实数据，相互印证；
  \item \textbf{多机制对照：}按 Abrams (2000) 思路比较恐惧进入方程的不同路径，避免单一通道偏误；
  \item \textbf{诚实边界：}对 $k$ 不可识别、Peacor 粒度、2D Turing 拓展未见显著斑图等局限明确说明。
\end{itemize}

\vspace{1em}
\noindent
\textbf{报告人：}李松茂、贺小轩、王玉麟 \hfill \textbf{日期：}2026 年 5 月 28 日

\end{document}
